\documentclass[11pt,a4paper]{article}

\usepackage[utf8]{inputenc}

\usepackage[T1]{fontenc}

\usepackage[polish]{babel}

\usepackage{lmodern}

\usepackage{geometry}

\geometry{margin=27mm}

\usepackage[hidelinks]{hyperref}

\title{Selektor jakości gaze przy zmianie telefonu i warunków}

\author{Mateusz Pomianek}

\date{Szkic naukowy, wersja 1.0.0 | 28 listopada 2026}

\begin{document}

\maketitle

\begin{abstract}

Proponowana miara główna to różnica ryzyka błędu wśród zachowanych klatek przy wspólnym pokryciu.

\end{abstract}

\section{Wprowadzenie}

Estymacja spojrzenia z kamer urządzeń mobilnych jest wykonalna w warunkach opisanych przez Krafkę et al. (2016). Kalibracja prawdopodobieństwa wymaga niezależnych danych i jawnej definicji trafności (Guo et al., 2017).

\section{Cel i pytania badawcze}

Hipotezy robocze wymagają zatwierdzenia przed rozpoczęciem właściwej weryfikacji.

\paragraph{H1.} Przy porównywalnym pokryciu wynoszącym 70 procent selektor oparty na jakości obrazu i geometrii twarzy osiągnie niższy odsetek błędnych zaakceptowanych pomiarów niż selektor oparty wyłącznie na pewności modelu bazowego.

\section{Materiał i metody}

Pierwszorzędowy wskaźnik: różnica ryzyka błędu wśród zachowanych klatek przy wspólnym pokryciu.

Trafność oznacza błąd kątowy nie większy niż 3 stopnie względem referencji, a pokrycie oznacza udział zachowanych klatek. Operacyjny próg ustali walidacja rozwojowa, a analiza krzywej ryzyko–pokrycie wykorzysta osobny telefon i nowych użytkowników. Wynik Briera opisze kalibrację.

Przed badaniem z udziałem ludzi należy zatwierdzić wymagania etyczne, zasady zgody i ochronę danych.

\section{Plan analizy}

Zamrożenie estymandu, reguł oceny jakości i obsługi braków nastąpi przed dostępem do danych wynikowych. Raport obejmie oszacowanie kontrastu głównego z przedziałem niepewności i pełne mianowniki. Jeśli przewidziano hipotezy drugorzędowe, przed analizą należy ustalić sposób kontroli wielokrotności.

\section{Wyniki}

W niniejszym pliku nie umieszczono wyników empirycznych, statystyk ani nieprzeprowadzonych analiz. Sekcja wymaga uzupełnienia po uzyskaniu i weryfikacji materiału dowodowego.

\section{Dyskusja i ograniczenia}

Ograniczenia obejmują możliwe błędy pomiaru, dobór próby, alternatywne mechanizmy i ograniczoną możliwość uogólniania. Korelacja nie dowodzi przyczynowości; poprawność techniczna nie jest dowodem znaczenia klinicznego.

\section{Wnioski}

Hipotezy i sposób pomiaru stanowią podstawę oceny wykonalności oraz przygotowania właściwego protokołu.

\section*{Bibliografia}

{\setlength{\parindent}{0pt}
\setlength{\parskip}{5pt}
\hangindent=1.27cm
Guo, C., Pleiss, G., Sun, Y., \& Weinberger, K. Q. (2017). On calibration of modern neural networks. \textit{Proceedings of the 34th International Conference on Machine Learning, 70}, 1321--1330. \url{https://proceedings.mlr.press/v70/guo17a.html}\par
\hangindent=1.27cm
Krafka, K., Khosla, A., Kellnhofer, P., Kannan, H., Bhandarkar, S., Matusik, W., \& Torralba, A. (2016). Eye tracking for everyone. \textit{Proceedings of the IEEE Conference on Computer Vision and Pattern Recognition}, 2176--2184. \url{https://openaccess.thecvf.com/content_cvpr_2016/html/Krafka_Eye_Tracking_for_CVPR_2016_paper.html}\par
}

\end{document}
